% ---------------------------------------------------------------------
\chapter{Elementare Datentypen}
\label{kap:01-datentypen}

\section*{Wochenziele}
\begin{wochenziele}
  \ziel{boolean, integer, float}{Mit den Grundtypen rechnen, Prioritäten von \texttt{not}, \texttt{and}, \texttt{or} beachten.}
  \ziel{Binärdarstellung}{integer und float ins Binärsystem umrechnen, periodische Brüche erkennen.}
  \ziel{IEEE 754}{float-Werte als 32-bit und 64-bit Fliesskommazahl angeben.}
  \ziel{set, list, matrix}{Mengenoperationen, komponentenweises Rechnen und Broadcasting anwenden.}
  \ziel{Funktion}{Den Begriff Funktion und Algorithmus erklären.}
\end{wochenziele}

\section{Kernideen}
Alle Daten im Computer sind Folgen von \textbf{Bits}. Daraus baut man die
Grundtypen: \textbf{boolean} (wahr/falsch), \textbf{integer} (ganze Zahlen) und
\textbf{float} (Kommazahlen mit endlich vielen Stellen). Mehrere Objekte fasst man
zu \textbf{Mengen} (ungeordnet), \textbf{Listen} (geordnet) und \textbf{Matrizen}
(zweidimensional) zusammen. Eine \textbf{Funktion} ordnet jedem Input genau einen
Output zu, ein \textbf{Algorithmus} beschreibt, wie man ihn berechnet.

% =====================================================================
\section{Bit und Boolean}

\begin{definition}{Bit}
  Die kleinste Einheit von Daten. Ein Bit kann genau \textbf{zwei} verschiedene
  Werte annehmen. Je nach Sichtweise:
  \texttt{1} und \texttt{0}, \texttt{Ein} und \texttt{Aus}, \texttt{Ja} und
  \texttt{Nein}, \texttt{Wahr} und \texttt{Falsch}.
\end{definition}

\begin{definition}{Boolescher Wert (boolean)}
  Ein boolescher Wert ist entweder \texttt{True} oder \texttt{False}.
  Ein Boolean ist also ein Bit, das man als \enquote{Wahr/Falsch} liest:
  \[
    \texttt{1} \;\widehat{=}\; \texttt{True} \;(T), \qquad
    \texttt{0} \;\widehat{=}\; \texttt{False} \;(F).
  \]
  Benannt nach George Boole (1815--1864), der als Erster 1 für True und 0 für
  False verwendet hat.
\end{definition}

\begin{beispiel}{Boolean in Python}
\begin{lstlisting}[language=Python]
True == 1        # True:  True verhält sich wie 1
False == 0       # True:  False verhält sich wie 0
True + True      # 2
bool(0)          # False: nur 0 ist False ...
bool(5)          # True:  ... jede andere Zahl ist True
\end{lstlisting}
\end{beispiel}

\subsection*{Einstellige und zweistellige Operatoren}
\begin{itemize}
  \item \textbf{Einstelliger Operator (unär):} Verknüpft \textbf{einen Operanden}
        und erzeugt daraus einen Wahrheitswert. Beispiel: $\neg T = F$ (\texttt{NOT}).
  \item \textbf{Zweistelliger Operator (binär):} Verknüpft \textbf{zwei Operanden}
        zu einem Ergebnis. Beispiel: $T \wedge T = T$ (\texttt{AND}).
\end{itemize}
\textbf{Fachlich entscheidend:} Die Einstelligkeit bzw.\ Zweistelligkeit
beschreibt die Anzahl der Operanden, auf die ein Operator angewendet wird.

\subsection*{Wahrheitstabelle}
\begin{center}
  \renewcommand{\arraystretch}{1.2}
  \begin{tabular}{cc|c|cc}
    \toprule
    $a$ & $b$ & \texttt{not} $a$ & $a$ \texttt{and} $b$ & $a$ \texttt{or} $b$ \\
    \midrule
    T & T & F & T & T \\
    T & F & F & F & T \\
    F & T & T & F & T \\
    F & F & T & F & F \\
    \bottomrule
  \end{tabular}
\end{center}

\begin{merke}
  \begin{itemize}[nosep]
    \item \texttt{not} dreht den Wert um.
    \item \texttt{and} ist \textbf{nur T}, wenn \textbf{beide} Operanden T sind.
    \item \texttt{or} ist \textbf{nur F}, wenn \textbf{beide} Operanden F sind.
  \end{itemize}
\end{merke}

\subsection*{Auswertungsreihenfolge}
\begin{enumerate}
  \item \textbf{Klammern} immer zuerst (innerste zuerst).
  \item \textbf{Priorität:} zuerst \texttt{not}, dann \texttt{and}, zuletzt \texttt{or}.
  \item Gleiche Operatoren von \textbf{links nach rechts} (Linksassoziativität).
  \item Ausnahme: \texttt{not not $x$} wird von \textbf{rechts} her gelesen
        (das innere \texttt{not} zuerst), also \texttt{not(not $x$)}.
\end{enumerate}

\begin{beispiel}{Serie 1, Aufgabe 1a}
  \texttt{False or True and False and True or False}\\
  $= F \text{ or } \underline{T \text{ and } F} \text{ and } T \text{ or } F$
  \quad (\texttt{and} zuerst, von links)\\
  $= F \text{ or } \underline{F \text{ and } T} \text{ or } F$\\
  $= \underline{F \text{ or } F} \text{ or } F = F \text{ or } F = \mathbf{F}$
\end{beispiel}

\begin{beispiel}{Serie 1, Aufgabe 1c}
  \texttt{not False or not not True and False or not True}\\
  $= T \text{ or } \underline{\text{not not } T} \text{ and } F \text{ or } F$
  \quad (\texttt{not} zuerst)\\
  $= T \text{ or } \underline{T \text{ and } F} \text{ or } F$\\
  $= T \text{ or } F \text{ or } F = \mathbf{T}$
\end{beispiel}

\begin{beispiel}{Serie 1, Aufgabe 1d (mit Klammern)}
  \texttt{not(not True or not(False and not True)) or False}\\
  $= \text{not}(F \text{ or not}(\underline{F \text{ and } F})) \text{ or } F$\\
  $= \text{not}(F \text{ or } \underline{\text{not } F}) \text{ or } F
   = \text{not}(\underline{F \text{ or } T}) \text{ or } F$\\
  $= \underline{\text{not } T} \text{ or } F = F \text{ or } F = \mathbf{F}$
\end{beispiel}

\begin{pruefung}
  Schritt für Schritt rechnen und in jeder Zeile nur \textbf{eine} Operation
  auswerten (unterstreichen). Die meisten Fehler entstehen, wenn man
  \texttt{or} vor \texttt{and} auswertet.
\end{pruefung}

% =====================================================================
\section{Integer}

\begin{definition}{Integer}
  Ein integer-Wert ist eine ganze Zahl:
  $\Z = \{\ldots, -2, -1, 0, 1, 2, \ldots\}$.
  Die Operatoren $-$ (Negation), $+$, $-$, $*$ sind immer wohldefiniert.
  Für die Division gibt es zwei eigene Operatoren:
  \begin{itemize}[nosep]
    \item \texttt{//} Ganzzahldivision: $10\,/\!/\,3 = 3$, \; $10\,/\!/\,7 = 1$, \; $10\,/\!/\,5 = 2$
    \item \texttt{\%} Rest: $10 \,\%\, 3 = 1$, \; $10 \,\%\, 7 = 3$, \; $10 \,\%\, 5 = 0$
  \end{itemize}
  Es gilt immer: $a = (a \,/\!/\, b)\cdot b + (a \,\%\, b)$.
\end{definition}

\begin{merke}
  In Python rundet \texttt{//} immer \textbf{nach unten} (Richtung $-\infty$):
  $-7\,/\!/\,2 = -4$ und $-7 \,\%\, 2 = 1$.
\end{merke}

\subsection*{Dezimal $\to$ binär (integer)}
\begin{enumerate}
  \item Zahl durch 2 teilen, Rest (0 oder 1) notieren.
  \item Mit dem Ergebnis von \texttt{//} weitermachen, bis 0 herauskommt.
  \item Reste \textbf{von unten nach oben} lesen.
\end{enumerate}

\begin{beispiel}{Serie 1, Aufgabe 2: $(409)_2$}
  \begin{tabular}{r@{\;}l}
    409 & Rest 1 \\ 204 & Rest 0 \\ 102 & Rest 0 \\ 51 & Rest 1 \\
    25 & Rest 1 \\ 12 & Rest 0 \\ 6 & Rest 0 \\ 3 & Rest 1 \\ 1 & Rest 1 \\ 0 &
  \end{tabular}
  \qquad $\Rightarrow\; (409)_2 = \mathbf{110011001}$ \;(von unten gelesen)\\[4pt]
  Kontrolle: $256 + 128 + 16 + 8 + 1 = 409$.
\end{beispiel}

\begin{pruefung}
  Immer zurückrechnen (Stellenwerte $1, 2, 4, 8, 16, \ldots$ addieren).
  Das kostet 20 Sekunden und findet fast jeden Fehler.
\end{pruefung}

% =====================================================================
\section{Float}

\subsection*{Periodische Dezimalzahl $\to$ Bruch}
Zahlen mit periodischen (oder endlichen) Nachkommastellen sind \textbf{rational}
($\in \mathbb{Q}$), also als Bruch darstellbar.
\begin{enumerate}
  \item Mit einer Zehnerpotenz multiplizieren, sodass die erste Periode
        \textbf{vor} dem Komma steht.
  \item Mit einer zweiten Zehnerpotenz multiplizieren, sodass die Periode
        \textbf{direkt nach} dem Komma beginnt.
  \item Subtrahieren: Die Periode fällt weg. Nach $x$ auflösen.
\end{enumerate}

\begin{beispiel}{Serie 1, Aufgabe 3: $2.10\overline{365}$}
  \[
    \begin{array}{r@{\;=\;}l}
      100000 \cdot x & 210365.\overline{365}\\
      -\quad 100 \cdot x & 210.\overline{365}\\ \hline
      99900 \cdot x & 210155
    \end{array}
    \quad\Rightarrow\quad x = \frac{210155}{99900} = \frac{42031}{19980}
  \]
\end{beispiel}

\subsection*{Nachkommastellen $\to$ binär}
\begin{enumerate}
  \item Nachkommateil mal 2 rechnen.
  \item Der Teil \textbf{vor} dem Komma (0 oder 1) ist das nächste Bit.
  \item Mit dem Teil nach dem Komma weitermachen.
  \item Stopp, wenn der Rest 0 ist, oder wenn ein Rest \textbf{wiederkommt}
        (dann ist die Darstellung periodisch).
  \item Bits \textbf{von oben nach unten} lesen.
\end{enumerate}

\begin{beispiel}{$0.625$ und $0.3$}
  \begin{minipage}[t]{0.45\linewidth}
    $0.625 \cdot 2 = \mathbf{1} + 0.25$\\
    $0.25 \cdot 2 = \mathbf{0} + 0.5$\\
    $0.5 \cdot 2 = \mathbf{1} + 0$\\[2pt]
    $\Rightarrow (0.625)_2 = 0.\mathbf{101}$
  \end{minipage}\hfill
  \begin{minipage}[t]{0.5\linewidth}
    $0.3 \cdot 2 = \mathbf{0} + 0.6$\\
    $0.6 \cdot 2 = \mathbf{1} + 0.2$ \;$\leftarrow$\\
    $0.2 \cdot 2 = \mathbf{0} + 0.4$\\
    $0.4 \cdot 2 = \mathbf{0} + 0.8$\\
    $0.8 \cdot 2 = \mathbf{1} + 0.6$ \;$\leftarrow$ \;0.6 kommt wieder\\[2pt]
    $\Rightarrow (0.3)_2 = 0.0\overline{1001}$
  \end{minipage}
\end{beispiel}

\begin{beispiel}{Serie 1, Aufgabe 4: $(1/6)_2$}
  $\tfrac16 \cdot 2 = \mathbf{0} + \tfrac13$, \quad
  $\tfrac13 \cdot 2 = \mathbf{0} + \tfrac23$, \quad
  $\tfrac23 \cdot 2 = \mathbf{1} + \tfrac13$ \;($\tfrac13$ kommt wieder)
  \quad$\Rightarrow\; (1/6)_2 = 0.0\overline{01}$
\end{beispiel}

\begin{merke}
  Irrationale Zahlen (z.\,B.\ $\pi$) und Zahlen mit unendlicher Periode
  (z.\,B.\ $0.3$ im Binärsystem) kann der Computer \textbf{nicht exakt}
  speichern. Ein \textbf{float} ist eine Kommazahl mit \textbf{endlich vielen}
  Nachkommastellen, es wird also gerundet. Darum ist in Python
  \texttt{0.7 + 0.1 == 0.9 - 0.1} \texttt{False}
  ($0.7999\ldots \ne 0.8$).
\end{merke}

\subsection*{Fliesskommazahl nach IEEE 754}
\[
  x = (-1)^{s} \cdot 1.m \cdot 2^{\,e - \text{Bias}}
\]
\begin{center}
  \renewcommand{\arraystretch}{1.2}
  \begin{tabular}{lcccc}
    \toprule
    Format & Vorzeichen $s$ & Exponent $e$ & Mantisse $m$ & Bias \\
    \midrule
    32-bit (\texttt{float} in Java)  & 1 bit & 8 bit  & 23 bit & 127 \\
    64-bit (\texttt{double} in Java) & 1 bit & 11 bit & 52 bit & 1023 \\
    \bottomrule
  \end{tabular}
\end{center}

\begin{enumerate}
  \item Zahl binär schreiben (Vorkomma- und Nachkommateil getrennt).
  \item Normalisieren: Komma verschieben bis $1.\ldots \cdot 2^{k}$.
  \item Vorzeichen: $s = 0$ für positiv, $s = 1$ für negativ.
  \item Exponent: $e = k + 127$ (bzw.\ $k + 1023$), binär mit 8 (bzw.\ 11) Bit,
        vorne mit Nullen auffüllen.
  \item Mantisse: die Bits \textbf{nach} der führenden 1 (die 1 wird nicht
        gespeichert). Auf 23 (bzw.\ 52) Bit kürzen und \textbf{runden}: Ist das
        erste abgeschnittene Bit eine 1 (und danach nicht nur Nullen), wird
        das letzte Bit um 1 erhöht.
\end{enumerate}

\begin{beispiel}{Serie 1, Aufgabe 5a: $0.00\overline{10}$ \;(= $1/6$)}
  $0.00101010\ldots = 1.\overline{01} \cdot 2^{-3}$\\[2pt]
  32-bit: $e = -3 + 127 = 124 = (01111100)_2$, Mantisse $0101\ldots01|\,1\ldots$
  wird aufgerundet:\\
  \hspace*{1em}\texttt{0 01111100 01010101010101010101011}\\[2pt]
  64-bit: $e = -3 + 1023 = 1020 = (01111111100)_2$, abgeschnittenes Bit ist 0,
  also abrunden:\\
  \hspace*{1em}\texttt{0 01111111100 0101\ldots0101} \;(52 Bit)
\end{beispiel}

\begin{beispiel}{Serie 1, Aufgabe 5b: $35.2$}
  $35 = (100011)_2$, \; $0.2 = (0.\overline{0011})_2$
  \;$\Rightarrow\; 35.2 = 100011.\overline{0011}_2 = 1.00011\overline{0011} \cdot 2^{5}$\\[2pt]
  32-bit: $e = 5 + 127 = 132 = (10000100)_2$, aufrunden:\\
  \hspace*{1em}\texttt{0 10000100 00011001100110011001101}\\[2pt]
  64-bit: $e = 5 + 1023 = 1028 = (10000000100)_2$, aufrunden:\\
  \hspace*{1em}\texttt{0 10000000100 00011001100\ldots110011010} \;(52 Bit)
\end{beispiel}

\begin{pruefung}
  \begin{itemize}[nosep]
    \item Bias auswendig: \textbf{127} (32-bit) und \textbf{1023} (64-bit).
    \item Die führende 1 kommt \textbf{nicht} in die Mantisse (\enquote{hidden bit}).
    \item Exponent immer auf volle Länge (8 bzw.\ 11 Bit) auffüllen.
    \item Bei periodischen Zahlen das \textbf{Runden} am Schluss nicht vergessen.
  \end{itemize}
\end{pruefung}

% =====================================================================
\section{Set, List und Matrix}

\begin{definition}{Menge (set)}
  Eine \textbf{ungeordnete} Sammlung von Objekten, \textbf{ohne} Rücksicht auf
  Wiederholung. $x \in M$: $x$ ist Element von $M$; $x \notin M$: ist es nicht.\\
  Aufzählen: $A = \{0, 1, \ldots, 9\}$ \quad Beschreiben:
  $B = \{x : x \text{ ist eine Ziffer}\}$
\end{definition}

\begin{center}
  \renewcommand{\arraystretch}{1.25}
  \begin{tabularx}{\linewidth}{@{}l l X@{}}
    \toprule
    Symbol & Name & Beispiel \\
    \midrule
    $=$ & Gleichheit & $\{a,b,c\} = \{b,c,a\}$, \; $\{a,b,c\} \ne \{a,b,x\}$ \\
    $\subset$ & Teilmenge & $\{a,b,c\} \subset \{a,b,c,x\}$, \; $\{\} \subset \{a,b,c\}$, \; $\{a,b\} \not\subset \{a,c,x\}$ \\
    $\cap$ & Durchschnitt & $\{a,b,c\} \cap \{b,c,d\} = \{b,c\}$, \; $\{a,b\} \cap \{c,d\} = \{\}$ \\
    $\cup$ & Vereinigung & $\{a,b,c\} \cup \{b,c,d\} = \{a,b,c,d\}$ \\
    $\setminus$ & Differenz & $\{a,b,c\} \setminus \{b,c,d\} = \{a\}$, \; $\{a,b\} \setminus \{c,d\} = \{a,b\}$ \\
    $\overline{\phantom{A}}$ & Komplement & in $G = \{a,b,c,d,e\}$: \; $\overline{\{a,b\}} = \{c,d,e\}$ \\
    \bottomrule
  \end{tabularx}
\end{center}

\begin{merke}
  \begin{itemize}[nosep]
    \item In dieser Vorlesung bedeutet $\subset$ \enquote{ist Teilmenge oder gleich}:
          $\{a,b,c\} \subset \{a,b,c\}$ ist wahr.
    \item Die leere Menge $\{\}$ ist Teilmenge \textbf{jeder} Menge.
    \item Das Komplement gibt es nur bezüglich einer Grundmenge $G$.
    \item Eine Menge mit $n$ Elementen hat $\mathbf{2^n}$ Teilmengen.
  \end{itemize}
\end{merke}

\begin{beispiel}{Serie 1, Aufgabe 6: Teilmengen von $\{a,b,c,d\}$ \;($2^4 = 16$)}
  \begin{tabular}{@{}lllll@{}}
    0 Elemente & 1 Element & 2 Elemente & 3 Elemente & 4 Elemente \\ \midrule
    $\{\}$ & $\{a\}, \{b\}$ & $\{a,b\}, \{a,c\}, \{a,d\}$ & $\{a,b,c\}, \{a,b,d\}$ & $\{a,b,c,d\}$ \\
           & $\{c\}, \{d\}$ & $\{b,c\}, \{b,d\}, \{c,d\}$ & $\{a,c,d\}, \{b,c,d\}$ & \\
    \midrule
    1 & 4 & 6 & 4 & 1
  \end{tabular}
\end{beispiel}

\begin{definition}{Liste (list)}
  Eine \textbf{geordnete} Sammlung von Objekten, \textbf{mit} Rücksicht auf
  Wiederholung, z.\,B.\ $a = (5, 3, 8, 8, 1)$. Die Objekte heissen Komponenten
  oder Einträge. $a_i$ bzw.\ \texttt{a[i]} ist die Komponente am $i$-ten Platz:
  $a_2 = \texttt{a[2]} = 3$ (Index ab 1).
\end{definition}

\begin{definition}{Matrix (matrix)}
  Eine zweidimensionale, rechteckig geordnete Sammlung mit Rücksicht auf
  Wiederholung. Eine $m \times n$ Matrix hat $m$ \textbf{Zeilen} und $n$
  \textbf{Spalten} (Dimension). $a_{ij} = \texttt{A[i,j]} = \texttt{A[i][j]}$
  ist der Eintrag in Zeile $i$, Spalte $j$.
  \[
    A = \begin{pmatrix} 5 & 3 & 8 & 8 & 1 \\ 1 & 0 & 0 & 9 & 7 \end{pmatrix}
    \quad (2 \times 5), \qquad a_{24} = 9
  \]
\end{definition}

\subsection*{Komponentenweise rechnen und Broadcasting}
\begin{itemize}
  \item \textbf{Komponentenweise:} Gleich grosse Listen/Matrizen werden Eintrag
        für Eintrag verrechnet.\\
        $(1,3,-1,0,1.5) + (1,2,2,1,0.5) = (2,5,1,1,2)$\\
        $(T,T,F) \text{ and } (F,T,F) = (F,T,F)$
  \item \textbf{Broadcasting:} Ein einzelner Wert wird zu einer Liste gleicher
        Grösse erweitert:\\
        $3 \cdot (2,5,4,-2) = (3,3,3,3) \cdot (2,5,4,-2) = (6,15,12,-6)$\\
        $3 - (0,1,2,3,2,1,0) = (3,2,1,0,1,2,3)$
  \item \textbf{Einstellige Operatoren} wirken auf alle Komponenten:
        $\text{not}(T,F,F,T,F) = (F,T,T,F,T)$
  \item Eine $m \times n$ Matrix lässt sich mit einer $1 \times 1$,
        $m \times 1$ (Spalte wird kopiert) oder $1 \times n$ (Zeile wird
        kopiert) Matrix verrechnen.
\end{itemize}

\begin{beispiel}{Broadcasting mit Matrizen (Serie 1, S.\,11)}
  \[
    \underbrace{2}_{1\times1} +
    \underbrace{\begin{pmatrix}1&2&1&1\\3&4&2&2\\3&5&1&1\end{pmatrix}}_{3\times4} -
    \underbrace{\begin{pmatrix}1\\2\\3\end{pmatrix}}_{3\times1} +
    \underbrace{\begin{pmatrix}0&1&1&2\end{pmatrix}}_{1\times4} \cdot
    \begin{pmatrix}1&2&3&2\\1&5&3&3\\1&2&2&5\end{pmatrix}
    = \begin{pmatrix}2&5&5&6\\3&9&5&8\\2&6&2&10\end{pmatrix}
  \]
  Hier ist $\cdot$ \textbf{komponentenweise} gemeint (nicht Matrixmultiplikation).
  Punkt vor Strich gilt weiterhin.
\end{beispiel}

\begin{pruefung}
  Vor dem Rechnen die Dimension jedes Teils notieren (wie oben unter den
  Klammern). So sieht man sofort, in welche Richtung kopiert wird und ob
  Broadcasting überhaupt möglich ist.
\end{pruefung}

% =====================================================================
\section{Funktion}

\begin{definition}{Funktion}
  Sei $D$ der Definitionsbereich (Inputs) und $W$ der Wertevorrat (Outputs).
  Eine Funktion $f\colon D \longrightarrow W$ ordnet \textbf{jedem} Element
  in $D$ \textbf{genau ein} Element in $W$ zu. Schreibweise:
  $x \mapsto y$ oder $f(x) = y$.
\end{definition}

\begin{definition}{Algorithmus}
  Eine präzise, schrittweise Angabe, die für jeden Input $x \in D$ beschreibt,
  wie der zugehörige Output $y = f(x)$ bestimmt wird.
\end{definition}

\begin{beispiel}{$\text{decToBin}\colon \N \to \{\text{Zeichenketten aus 0 und 1}\}$}
\begin{lstlisting}[language=Python]
def decToBin(x):
    y = ''
    temp = x
    while temp != 0:
        if temp % 2 == 0:
            y = '0' + y      # Rest 0 vorne anhängen
        else:
            y = '1' + y      # Rest 1 vorne anhängen
        temp = temp // 2
    return y
\end{lstlisting}
  Das ist genau das Verfahren \enquote{durch 2 teilen, Reste von unten lesen}:
  Jeder neue Rest wird \textbf{vorne} angehängt.
\end{beispiel}

% =====================================================================
\section{Prüfungsrelevant}
\begin{pruefung}
  \begin{itemize}[nosep]
    \item Boolesche Ausdrücke: Klammern, dann \texttt{not} $>$ \texttt{and} $>$ \texttt{or}, sonst von links.
    \item \texttt{and} nur T bei T/T, \texttt{or} nur F bei F/F.
    \item Integer $\to$ binär: durch 2 teilen, Reste \textbf{von unten} lesen.
    \item Nachkomma $\to$ binär: mal 2, Bits \textbf{von oben} lesen, auf Wiederholung achten.
    \item Periode $\to$ Bruch: zwei Gleichungen, subtrahieren.
    \item IEEE 754: $1|8|23$ mit Bias 127, $1|11|52$ mit Bias 1023, hidden bit, runden.
    \item Menge: ungeordnet, ohne Wiederholung. Liste: geordnet, mit Wiederholung.
    \item $n$ Elemente $\Rightarrow 2^n$ Teilmengen.
    \item Broadcasting: $1\times1$, $m\times1$, $1\times n$ mit $m\times n$.
    \item Funktion: jedem Input \textbf{genau ein} Output.
  \end{itemize}
\end{pruefung}

% =====================================================================
\section{Kontrollfragen}
\label{sec:01-fragen}
Erst selbst lösen, dann in Abschnitt~\ref{sec:01-loesungen} vergleichen.

\subsection*{Bit und Boolean}
\begin{enumerate}[label=\textbf{F\arabic*}, ref=F\arabic*, series=fragen01]
  \item Was ist ein Bit? Nenne vier Sichtweisen auf seine zwei Werte.
  \item Was ist der Unterschied zwischen einem einstelligen und einem
        zweistelligen Operator? Nenne je ein Beispiel.
  \item Werte aus: \texttt{True and not False or False}
  \item Werte aus: \texttt{False or not False and False}
  \item Werte aus: \texttt{not(True and (False or not False)) or not True}
  \item Welchen Wert liefern in Python \texttt{True + True} und \texttt{bool(0)}?
\end{enumerate}

\subsection*{Integer}
\begin{enumerate}[resume*=fragen01]
  \item Berechne $23\,/\!/\,4$, $23 \,\%\, 4$ und in Python $-7\,/\!/\,2$.
  \item Bestimme $(100)_2$.
  \item Welche Dezimalzahl ist $(101101)_2$?
\end{enumerate}

\subsection*{Float}
\begin{enumerate}[resume*=fragen01]
  \item Schreibe $1.2\overline{3}$ als Bruch.
  \item Bestimme $(0.375)_2$ und $(0.1)_2$.
  \item Warum liefert \texttt{0.7 + 0.1 == 0.9 - 0.1} in Python \texttt{False}?
  \item Gib $-6.5$ als 32-bit Fliesskommazahl an.
  \item Warum wird die führende 1 der Mantisse nicht gespeichert?
\end{enumerate}

\subsection*{Set, List, Matrix}
\begin{enumerate}[resume*=fragen01]
  \item Wie viele Teilmengen hat $\{1,2,3,4,5\}$?
  \item Sei $G = \{1,\ldots,6\}$, $A = \{1,2,3,4\}$, $B = \{3,4,5\}$.
        Bestimme $A \cap B$, $A \cup B$, $A \setminus B$, $B \setminus A$ und $\overline{A}$.
  \item Ist $\{1,1,2\} = \{1,2\}$? Ist $(1,1,2) = (1,2)$? Ist $(1,2) = (2,1)$?
  \item Berechne $2 \cdot (1,-1,3) - 1$.
  \item Mit welchen Dimensionen kann eine $3 \times 4$ Matrix per Broadcasting
        verrechnet werden? Geht $4 \times 1$?
\end{enumerate}

\subsection*{Funktion}
\begin{enumerate}[resume*=fragen01]
  \item Ist die Zuordnung $x \mapsto \pm\sqrt{x}$ (für $x > 0$) eine Funktion?
  \item Was liefert \texttt{decToBin(6)}? Und \texttt{decToBin(0)}?
\end{enumerate}

% =====================================================================
\section{Lösungen zu den Kontrollfragen}
\label{sec:01-loesungen}
\begin{enumerate}[label=\textbf{F\arabic*}]
  \item Kleinste Dateneinheit mit genau zwei Werten:
        1/0, Ein/Aus, Ja/Nein, Wahr/Falsch.
  \item Einstellig: ein Operand (\texttt{not}, Negation $-x$).
        Zweistellig: zwei Operanden (\texttt{and}, \texttt{or}, $+$, $\cap$).
  \item $T \text{ and } \underline{\text{not } F} \text{ or } F
         = \underline{T \text{ and } T} \text{ or } F = T \text{ or } F = \mathbf{T}$
  \item $F \text{ or } \underline{\text{not } F} \text{ and } F
         = F \text{ or } \underline{T \text{ and } F} = F \text{ or } F = \mathbf{F}$
  \item $\text{not}(T \text{ and } (F \text{ or } T)) \text{ or } F
         = \text{not}(T \text{ and } T) \text{ or } F = F \text{ or } F = \mathbf{F}$
  \item \texttt{True + True} $= 2$ (True zählt als 1), \texttt{bool(0)} $=$ \texttt{False}.
  \item $23\,/\!/\,4 = 5$, \; $23 \,\%\, 4 = 3$ (Kontrolle: $5 \cdot 4 + 3 = 23$),
        \; $-7\,/\!/\,2 = -4$ (abrunden).
  \item $100, 50, 25, 12, 6, 3, 1 \to$ Reste $0,0,1,0,0,1,1$
        $\Rightarrow (100)_2 = 1100100$ \;($64 + 32 + 4$).
  \item $32 + 8 + 4 + 1 = 45$.
  \item $100x = 123.\overline{3}$, \; $10x = 12.\overline{3}$
        $\Rightarrow 90x = 111 \Rightarrow x = \frac{111}{90} = \frac{37}{30}$.
  \item $0.375 \to 0.75\,(0) \to 1.5\,(1) \to 1.0\,(1)$: $(0.375)_2 = 0.011$.\\
        $0.1 \to 0.2\,(0) \to 0.4\,(0) \to 0.8\,(0) \to 1.6\,(1) \to 1.2\,(1) \to 0.4 \ldots$:
        $(0.1)_2 = 0.0\overline{0011}$.
  \item $0.1$, $0.7$, $0.9$ sind binär periodisch, also nicht exakt speicherbar.
        Durch das Runden ist $0.7 + 0.1 = 0.7999\ldots$, aber $0.9 - 0.1 = 0.8$.
  \item $6.5 = 110.1_2 = 1.101 \cdot 2^2$; \; $s = 1$, $e = 2 + 127 = 129 = 10000001_2$:\\
        \texttt{1 10000001 10100000000000000000000}
  \item Nach dem Normalisieren steht vor dem Komma \textbf{immer} eine 1.
        Sie muss nicht gespeichert werden, so gewinnt man ein Bit Genauigkeit.
  \item $2^5 = 32$.
  \item $A \cap B = \{3,4\}$, \; $A \cup B = \{1,2,3,4,5\}$, \;
        $A \setminus B = \{1,2\}$, \; $B \setminus A = \{5\}$, \;
        $\overline{A} = \{5,6\}$.
  \item Ja (Mengen ignorieren Wiederholung). Nein (Liste zählt Wiederholung).
        Nein (Liste ist geordnet).
  \item $(2,-2,6) - (1,1,1) = (1,-3,5)$.
  \item $1 \times 1$, $3 \times 1$ und $1 \times 4$. \; $4 \times 1$ geht
        \textbf{nicht}, weil die Zeilenzahl (4 statt 3) nicht passt.
  \item Nein: Zu jedem $x > 0$ gibt es \textbf{zwei} Outputs, eine Funktion
        verlangt genau einen.
  \item \texttt{'110'}. Für \texttt{0} läuft die Schleife nie, also kommt die
        leere Zeichenkette \texttt{''} zurück (richtig wäre \texttt{'0'}).
\end{enumerate}

\section{Offene Fragen}
% Was ist noch unklar? Beim Dozenten oder in der Übung nachfragen.
